%% kapa: exam=1 round=24 year=2013 period=2 form=A subject=appraisal_law subject_ko=부동산관계법규 seq=11 number=51
%% source: Q-Net 감정평가사 기출문제 (cst00309) artlSeq=5073802 file=extracted/1008/5073802/2013년 제24회 감정평가사 1차시험 기출문제/2교시(A형).hwp
%% conversion: kordoc 4.24.51 -> tools/gen1.py (auto); review=proofread-sonnet; flags=none
\begin{question}{51}
국토의 계획 및 이용에 관한 법령상 개발행위허가에 관한 설명으로 옳은 것은?
\begin{choices}
  \item 개발행위허가를 받은 행정청이 새로 공공시설을 설치한 경우 새로 설치된 공공시설은 그 시설을 관리할 관리청에 유상으로 귀속된다.
  \item 개발밀도관리구역 안에서 개발행위허가 신청서를 제출할 때에도 기반시설의 설치 계획서를 첨부하여야 한다.
  \item 허가권자는 개발행위 불허가처분을 할 경우 불허가처분의 사유를 지체 없이 구두로 알려야 한다.
  \item 개발행위허가에는 조건을 붙일 수 없다.
  \item 허가권자는 개발행위허가의 신청 내용이 도시ㆍ군관리계획의 내용에 어긋나는 경우 개발행위허가를 하여서는 아니된다.
\end{choices}
\end{question}
\answer{5}
